\section{Conclusion}
\label{sec:conclusion}

Natural language to SQL has long been framed as a translation problem. SQLclMCP reframes it as an \emph{end-to-end verification problem}: generating a query is only half the challenge; confirming that it is semantically correct, executes efficiently, and preserves optimizer behavior is the other half. By tackling all three dimensions simultaneously on Oracle Database 26ai, this work establishes a new standard for rigorous, production-oriented NL2SQL evaluation.

\subsection{What Makes This Framework Different}

Most NL2SQL benchmarks stop at exact-match or execution-accuracy. SQLclMCP goes further by introducing a \textbf{three-layer validation stack} that no prior public evaluation combines:

\begin{description}
    \item[Layer 1 — Semantic Correctness.]
    Normalized, sorted result-set comparison confirms that generated SQL returns the \emph{right answer}, not just syntactically valid SQL. On 120 TPC-H questions spanning three complexity tiers, this layer records \textbf{99.17\% accuracy} (119/120) — 100\% on simple and medium queries, 97.73\% on complex queries.

    \item[Layer 2 — Execution Performance.]
    Wall-clock latency is measured for both baseline and MCP-generated queries across every complexity tier. The framework achieves a \textbf{2.8\% reduction in mean latency} and a striking \textbf{30.3\% reduction in P95 latency}, demonstrating that the model consistently avoids pathological query shapes that inflate tail execution times.

    \item[Layer 3 — Optimizer Parity via EXPLAIN PLAN.]
    This is the most distinctive contribution. By comparing Oracle EXPLAIN PLAN output at the cost and cardinality level, SQLclMCP verifies that generated queries do not merely return correct answers — they follow \emph{the same execution path} as hand-written SQL. \textbf{100\% cost agreement and 100\% cardinality agreement} across all 120 queries confirm that the Oracle query optimizer treats MCP-generated SQL as first-class, plan-stable input. No prior NL2SQL evaluation has demonstrated this property on a commercial enterprise database.
\end{description}

Together, these three layers transform evaluation from a pass/fail accuracy report into a \emph{production readiness certificate}. An enterprise deploying SQLclMCP can answer not just ``does it return the right rows?'' but also ``will it stress my indexing strategy?'', ``will it cause latency spikes under load?'', and ``does the optimizer understand it?''

\subsection{Key Results at a Glance}

\begin{itemize}
    \item \textbf{99.17\% semantic accuracy} on 120 leakage-safe TPC-H queries (Oracle 26ai)
    \item \textbf{100\% EXPLAIN PLAN parity} on both cost and cardinality estimation
    \item \textbf{30.3\% P95 latency reduction} — the MCP model avoids worst-case query plans
    \item \textbf{Single, deterministically correctable failure}: MySQL \texttt{YEAR()} used instead of Oracle \texttt{EXTRACT(YEAR FROM ...)} — a dialect mismatch patched by schema-aware post-processing
    \item Complexity-stratified results with full reproducibility via open JSON artifacts
\end{itemize}

\subsection{Broader Significance}

The implications reach beyond benchmark numbers. Three forces converge to make this moment pivotal:

\textbf{Democratization at scale.} At 99\%+ accuracy, business analysts, data scientists, and domain experts can query Oracle databases in plain English without writing a single line of SQL. The productivity gains — fewer hand-offs to database engineering teams, faster iteration on analytical questions, lower error rates from mis-typed joins — are compounding.

\textbf{Trust through transparency.} EXPLAIN PLAN parity is not just a technical metric; it is an audibility guarantee. Operations teams can inspect generated query plans just as they would hand-written SQL. This lowers the organizational barrier to adopting AI-generated queries in regulated industries where ``black-box'' outputs are unacceptable.

\textbf{The MCP protocol as an evaluation interface.} By routing evaluation traffic through the Model Context Protocol, SQLclMCP establishes MCP as a viable standard interface for benchmarking LLM database tools — a replicable design that other researchers and teams can adopt for any LLM-database pair.

\subsection{Limitations and Honest Caveats}

These results come with important caveats:
\begin{itemize}
    \item Evaluation targets Oracle 26ai exclusively; cross-database generalization (PostgreSQL, MySQL, MSSQL) requires separate study.
    \item The test set comprises 120 TPC-H analytical queries; OLTP workloads and schemas with 100+ tables are not covered.
    \item A single LLM backend (Claude API) is evaluated; comparative benchmarking against GPT-4, Gemini, or open-source alternatives remains open work.
    \item Statistical power is moderate at $n=120$; a 500+ query corpus would sharpen confidence intervals on failure rate estimates.
\end{itemize}

\subsection{Future Directions}

The framework itself is designed to be extensible. Near-term priorities include expanding the query corpus, adding multi-database targets, and integrating schema-aware retry logic to eliminate the remaining dialect-mismatch failures. Medium-term, EXPLAIN PLAN comparisons can be elevated to \emph{cost-budget enforcement} — automatically rejecting generated queries whose estimated cost exceeds a configurable threshold, enabling fine-grained SLA compliance for analytics workloads.

Longer-term, the combination of semantic correctness signals and optimizer feedback creates an ideal training loop for fine-tuning models that are not only accurate but \emph{plan-aware}: models that learn to prefer index-friendly access patterns, avoid Cartesian products, and select join orders that preserve cardinality estimates.

\subsection{Final Remarks}

SQLclMCP demonstrates that the NL2SQL challenge is no longer whether language models can generate correct SQL — at 99.17\%, that question is largely answered for analytical workloads. The frontier has shifted to \emph{depth of validation}: can we certify not only that the answer is correct, but that the path to it is efficient, plan-stable, and production-safe?

With three-layer validation anchored by EXPLAIN PLAN optimizer parity, this framework answers that question affirmatively — and opens an exciting path toward AI-powered SQL generation that enterprises can deploy with confidence, audit without reservation, and trust at scale.

\subsection{Reproducibility}

All evaluation artifacts are publicly archived for reproducibility:
\begin{itemize}
    \item Source code: GitHub repository (SQLclMCP)
    \item Test questions: 120-question dataset in JSON format with complexity labels
    \item Results data: Raw JSON outputs with per-query metrics
    \item Evaluation code: Python framework with Oracle and Claude API integration
    \item Configuration: Docker Compose for one-command reproducibility
\end{itemize}

Researchers can replicate evaluation, extend to additional database systems, and validate results on their own schemas.
